\section*{Chapter \ref{ch:m1:clearing-and-settlement} --- Clearing and Settlement}

\begin{solution}{exo:m1:clearing-and-settlement:1}
$40\times39/2 = 780$ bilateral relationships against 40: a factor of 19.5,
that is $(n-1)/2$.
\end{solution}

\begin{solution}{exo:m1:clearing-and-settlement:2}
A: $-2\,000 + 1\,500 = -500$ shares, cash $+100\,000 - 75\,150 =
+\$24\,850$. B: $0$ shares, $-100\,000 + 100\,400 = +\$400$. C: $+500$
shares, $-100\,400 + 75\,150 = -\$25\,250$. Gross value \$275\,550; net cash
moving \$25\,250: efficiency 90.8\%.
\end{solution}

\begin{solution}{exo:m1:clearing-and-settlement:3}
Two days: $2.33 \times 0.025 \times \sqrt2 \times 400 = \$33.0$ million. One
day: \$23.3 million. The shorter cycle frees \$9.7 million, 29\% of the
requirement.
\end{solution}

\begin{solution}{exo:m1:clearing-and-settlement:4}
(a) 120 from the margin, 20 from the defaulter's fund share; nobody else
pays. (b) $120 + 30 + 20 + 130$ from the survivors' fund: \$5.2 million per
surviving member. (c) $120 + 30 + 20 + 400 + 400$, and \$30 million is
uncovered: the \omterm{def:m1:clearing-and-settlement:clearing}{clearing} house is insolvent unless its rules allow further
loss allocation.
\end{solution}

\begin{solution}{exo:m1:clearing-and-settlement:5}
Before: $2.33\times0.03\times\sqrt2\times600 = \$59$ million. After: \$356
million, six times more, for an unchanged position. The new requirement is
71\% of the broker's capital: it would be reasonable for the \omterm{def:m1:clearing-and-settlement:clearing}{clearing} house
to doubt that the broker can fund a further rise, which is what a capital
premium expresses, and rational for the broker to stop the position from
growing.
\end{solution}

\begin{solution}{exo:m1:clearing-and-settlement:6}
Without DVP the seller has delivered and receives nothing: it loses \$10
million and becomes a creditor in a bankruptcy. With DVP it still has its
shares, now worth \$9.6 million: a replacement loss of \$400\,000. With a
CCP it loses nothing: the CCP performs at the original price and bears the
\$400\,000 against the buyer's margin.
\end{solution}

\begin{solution}{exo:m1:clearing-and-settlement:7}
96.0\% before; 94.3\% after (the exact second figure depends on the random
counterparties chosen). The 400 added trades are all in one direction for one
member in one stock: nothing offsets them, so they add \$4 million of gross
value of which almost all must settle, and the ratio falls. One-directional
flow is what netting cannot help --- and what margin is charged on.
\end{solution}

\begin{solution}{exo:m1:clearing-and-settlement:8}
Statistical: defaults of large members are rare events; ``no use in $N$
years'' is a sample with almost no observations of the event the fund exists
for, and says nothing about a loss sized to the tail. Margins at 99\% are
\emph{designed} to be exceeded one period in a hundred per member.
Behavioural: the fund is also what makes members monitor one another and the
\omterm{def:m1:clearing-and-settlement:clearing}{clearing} house's admission and margin policy; shrink it and both the
resources and the discipline fall, and margins were calibrated in the
presence of that discipline. The relevant test is a stress scenario (the
default of the two largest members in extreme conditions), not history.
\end{solution}

\begin{solution}{pb:m1:clearing-and-settlement:1}
\textbf{1.} $2.33 \times 0.013 \times \sqrt2 \times 2\,400 = \$102.8$
million.
\textbf{2.} $z = 130/(0.013\sqrt2 \times 2\,400) = 2.95$: 99.8\%.
\textbf{3.} $0.03 \times 2\,400 = \$72$ million.
\textbf{4.} Zero: the loss has been paid, and the position, now worth
\$2\,328 million, is again covered by \$130 million.
\textbf{5.} $0.06 \times 2\,328 = \$139.7$ million.
\textbf{6.} The position is worth \$2\,188 million at the open; a further 2\%
costs \$43.8 million: \$183.4 million in total.
\textbf{7.} $0.004 \times 2\,144.6 = \$8.6$ million; total \$192.0 million.
\textbf{8.} 130 (C's margin), 25 (C's fund contribution), 15 (\omterm{def:m1:clearing-and-settlement:clearing}{clearing}
house capital), 22.0 (survivors' fund).
\textbf{9.} $22.0/20 = \$1.10$ million each.
\textbf{10.} Loss $139.7 + 0.004 \times 2\,188.3 = \$148.4$ million, less
than C's own \$155 million: neither the \omterm{def:m1:clearing-and-settlement:clearing}{clearing} house's capital nor the
survivors would have been touched. Three hours cost them \$37 million.
\textbf{11.} \$192 million, that is 6.2 times $\sigma N$ --- against 2.95
held.
\textbf{12.} $130+25+15+400+400 = \$970$ million.
\textbf{13.} Solving $2\,328 f + 0.004 \times 2\,328(1-f) = 970$ gives $f =
41\%$.
\textbf{14.} $3 \times 0.013 \times \sqrt3 \times 2\,400 = \$162$ million.
\textbf{15.} A charge increasing in the ratio of the position to the market's
capacity to absorb it, for instance proportional to notional times
$\sqrt{\text{position}/\text{daily volume}}$: the cost of liquidating grows
like the square root of size (\cref{met:m1:the-sell-side:sqrt}), and a member
holding 30\% of the open interest \emph{is} the market it would be liquidated
into.
\textbf{16.} So that the \omterm{def:m1:clearing-and-settlement:clearing}{clearing} house loses its own money before its
members lose theirs: those who set the margin model must suffer from its
errors first.
\textbf{17.} Unlimited assessments make every member's liability to the
\omterm{def:m1:clearing-and-settlement:clearing}{clearing} house unbounded and unquantifiable: their own capital requirements
and risk limits cannot be set, the best-capitalised members leave, and in a
crisis the assessments would propagate the default instead of containing
it.
\textbf{18.} Margin calls rise for all members when volatility rises, draining
cash exactly when it is scarcest and forcing sales that raise volatility
further: procyclicality.
\textbf{19.} \textbf{\$22.0 million.}
\textbf{20.} The other members, jointly, up to the limit of what they have
committed --- and beyond that limit, nobody.
\end{solution}

\begin{solution}{iq:m1:clearing-and-settlement:1}
\omterm{def:m1:clearing-and-settlement:ccp}{Novation} replaces the contract between the two trading parties by two
contracts with the \omterm{def:m1:clearing-and-settlement:clearing}{clearing} house in the middle. A firm no longer bears the
credit risk of whoever it happened to trade with on an anonymous order book,
only that of the \omterm{def:m1:clearing-and-settlement:clearing}{clearing} house; and its positions against all
counterparties net into one.
\iqlookfor{counterparty risk and netting, both.}
\end{solution}

\begin{solution}{iq:m1:clearing-and-settlement:2}
\omterm{def:m1:clearing-and-settlement:margin}{Variation margin} is the daily \omterm{def:m1:clearing-and-settlement:clearing}{settlement} of gains and losses in cash: it is
paid away to whoever won, and is not returned. \omterm{def:m1:clearing-and-settlement:margin}{Initial margin} is collateral
against future loss during a close-out: it remains the member's property and
comes back when the position is closed.
\iqlookfor{a payment versus a deposit.}
\end{solution}

\begin{solution}{iq:m1:clearing-and-settlement:3}
Customers' purchases are guaranteed by the \omterm{def:m1:clearing-and-settlement:clearing}{clearing} house from trade date to
\omterm{def:m1:clearing-and-settlement:clearing}{settlement}, then two days later; the broker posts margin for them,
proportional to net unsettled positions and to volatility. One-directional
buying in stocks whose volatility had exploded multiplied the requirement in
days, beyond the broker's liquid capital. Blocking purchases (not sales) was
the one action that stopped the requirement growing; the alternative was
failing the margin call, that is, defaulting.
\iqlookfor{margin $\propto$ net position $\times$ volatility $\times
\sqrt{\text{days}}$, and that sales reduce it.}
\end{solution}

\begin{solution}{iq:m1:clearing-and-settlement:4}
Defaulter's margin; defaulter's default-fund contribution; a tranche of the
\omterm{def:m1:clearing-and-settlement:clearing}{clearing} house's capital; survivors' default-fund contributions; capped
assessments. Defaulter-pays first so that each member internalises its own
risk; the \omterm{def:m1:clearing-and-settlement:clearing}{clearing} house next so that it has an incentive to margin
properly; mutualised layers last, to cover the tail while giving members a
reason to police admission and concentration.
\iqlookfor{the incentive logic of the order, not just the list.}
\end{solution}

\begin{solution}{iq:m1:clearing-and-settlement:5}
Assert: net shares per symbol sum to zero across members; net cash per
currency sums to zero; every accepted trade appears in exactly two members'
instructions; counts and gross totals match the matching engine's
end-of-day figures. If one fails: do \emph{not} publish; escalate at once to
operations and the \omterm{def:m1:clearing-and-settlement:clearing}{settlement} agent to obtain the extension that exists for
this case; find the missing or duplicated trade from the reconciliation
difference. Publishing wrong instructions is far worse than publishing late.
\iqlookfor{refusing to publish, and knowing that an escalation path must
exist before the day it is needed.}
\end{solution}

\begin{solution}{iq:m1:clearing-and-settlement:6}
For: multilateral netting shrinks exposures; margin is collected from
everyone under one transparent model; a default is managed by an organised
auction instead of thousands of bilateral disputes, as in September 2008;
supervisors see positions. Against: risk is concentrated in a few entities
too important to fail; margin is procyclical and transmits stress through
liquidity; members' mutual guarantees link all large banks; netting across
products is lost when \omterm{def:m1:clearing-and-settlement:clearing}{clearing} is split by asset class; and the products that
remain uncleared are the hardest ones.
\iqlookfor{concentration and procyclicality named on the ``against'' side.}
\end{solution}
